\documentclass{amsart}
\usepackage{amsmath,amssymb,amsthm,hyperref}
\theoremstyle{plain}
\newtheorem{theorem}{Theorem}
\newtheorem{lemma}{Lemma}
\newtheorem{proposition}{Proposition}
\newtheorem{corollary}{Corollary}
\theoremstyle{definition}
\newtheorem{definition}{Definition}
\newtheorem{remark}{Remark}
\newtheorem{example}{Example}

\newcommand{\FI}{\mathcal{FI}}
\newcommand{\Fcal}{\mathcal{F}}

\begin{document}

\title{Lattice closure of the exchange principle inside the standard families of fuzzy implications}
\maketitle

\section{Statement and notation}

Throughout, $\FI$ is the set of fuzzy implications on $[0,1]$: functions
$I\colon[0,1]^2\to[0,1]$ that are decreasing in the first variable, increasing in the
second variable, and satisfy
\begin{equation}\label{eq:bnd}
I(0,0)=I(1,1)=1,\qquad I(1,0)=0 .
\end{equation}
The pointwise lattice operations are
\[
(I\vee J)(x,y)=\max\{I(x,y),J(x,y)\},\qquad
(I\wedge J)(x,y)=\min\{I(x,y),J(x,y)\}.
\]
We write $I\le J$ when $I(x,y)\le J(x,y)$ for all $(x,y)\in[0,1]^2$, and we call $I,J$
\emph{comparable} when $I\le J$ or $J\le I$. We say $I$ satisfies the \emph{exchange
principle} (EP) if
\begin{equation}\label{eq:EP}
I(x,I(y,z))=I(y,I(x,z))\qquad\text{for all }x,y,z\in[0,1].
\end{equation}
We say $I$ satisfies the \emph{left neutrality property} (NP) if
\begin{equation}\label{eq:NP}
I(1,y)=y\qquad\text{for all }y\in[0,1].
\end{equation}

The four standard subfamilies are, as in the statement:
\begin{itemize}
\item $(S,N)$-implications $I_{S,N}(x,y)=S(N(x),y)$, $S$ a t-conorm, $N$ a fuzzy negation;
\item $R$-implications $I_T(x,y)=\sup\{t\in[0,1]:T(x,t)\le y\}$, $T$ a t-norm;
\item $f$-implications $I_f(x,y)=f^{-1}\!\bigl(x\,f(y)\bigr)$, $f\colon[0,1]\to[0,\infty]$
strictly decreasing and continuous with $f(1)=0$ (with the conventions
$0\cdot\infty=0$ and $f^{-1}(t)=0$ for $t\ge f(0)$);
\item $g$-implications $I_g(x,y)=g^{-1}\!\bigl(\min\{g(y)/x,\,g(1)\}\bigr)$,
$g\colon[0,1]\to[0,\infty]$ strictly increasing and continuous with $g(0)=0$
(with $g(y)/0:=\infty$).
\end{itemize}

\paragraph{Summary of results and a correction.}
A natural first guess is the clean dichotomy ``$I\vee J$ satisfies (EP) $\iff$ $I\wedge J$
satisfies (EP) $\iff$ $I,J$ comparable''.  This dichotomy is \emph{correct for the
$f$- and $g$-families} but \emph{false for the $(S,N)$-family}.  The decisive structural feature
is whether the implications of the family are \emph{strictly} monotone on the interior of
their range:
\begin{itemize}
\item $f$- and $g$-implications are strictly monotone in the interior, and there the
dichotomy holds (Theorem~\ref{thm:fg}).
\item $(S,N)$-implications need not be strictly monotone, and there is a genuine
\emph{merging} phenomenon: two $(S,N)$-implications built from the \emph{same} t-conorm
$S$ but \emph{incomparable} negations have a join (and a meet) that is again an
$(S,N)$-implication, hence satisfies (EP).  Thus comparability is \emph{not} necessary.
We give an explicit counterexample (Example~\ref{ex:SN}) and the corrected statement
(Theorem~\ref{thm:SN}).
\end{itemize}
The common engine of the positive results is the reduction
Lemma~\ref{lem:reduction}, valid for arbitrary implications.  The necessity direction for
the $f$- and $g$-families is established by a completely elementary algebraic computation
(Section~\ref{sec:fg}), with no appeal to numerics.

\section{Uniform properties of the four families}

\begin{lemma}[Neutrality and boundary values]\label{lem:NP}
Every member $I$ of each of the four families satisfies \textup{(NP)} and, in addition,
\begin{equation}\label{eq:topright}
I(x,1)=1\quad\text{and}\quad I(0,y)=1\qquad\text{for all }x,y\in[0,1].
\end{equation}
\end{lemma}

\begin{proof}
We check \eqref{eq:NP} for each family.

\emph{$(S,N)$.} A t-conorm satisfies $S(0,y)=y$, and a fuzzy negation satisfies $N(1)=0$;
hence $I_{S,N}(1,y)=S(N(1),y)=S(0,y)=y$.

\emph{$R$.} A t-norm satisfies $T(1,t)=t$, so
$I_T(1,y)=\sup\{t:T(1,t)\le y\}=\sup\{t:t\le y\}=y$.

\emph{$f$.} Since $f$ is a strictly decreasing continuous bijection from $[0,1]$ onto
$[0,f(0)]$ with $f(1)=0$, $f^{-1}$ is its (strictly decreasing) inverse and
$I_f(1,y)=f^{-1}(1\cdot f(y))=f^{-1}(f(y))=y$.

\emph{$g$.} Since $g(y)\le g(1)$ for all $y$, $\min\{g(y)/1,g(1)\}=g(y)$, hence
$I_g(1,y)=g^{-1}(g(y))=y$.

For \eqref{eq:topright}: $I(x,1)\ge I(x,y)$ for all $y$ because $I$ is increasing in the
second variable; in each family one checks $I(x,1)=1$ directly
($S(N(x),1)=1$; $\sup\{t:T(x,t)\le 1\}=1$; $f^{-1}(x f(1))=f^{-1}(0)=1$;
$g^{-1}(\min\{g(1)/x,g(1)\})=g^{-1}(g(1))=1$). Finally $I(0,y)\ge I(0,0)=1$ by
monotonicity in the first variable together with $I(0,0)=1$ from \eqref{eq:bnd}, so
$I(0,y)=1$.
\end{proof}

\begin{lemma}[Continuity in the generic case]\label{lem:cont}
For $f$- and $g$-implications the implication $I$ is automatically continuous on
$[0,1]^2$ from the stated hypotheses on $f,g$.  For $(S,N)$- and $R$-implications $I$ is
continuous when $S$ and $N$, respectively the t-norm $T$, is continuous.
\end{lemma}

\begin{proof}
For $f$-implications, $f$ and $f^{-1}$ are continuous and $(x,y)\mapsto x f(y)$ is
continuous into $[0,\infty]$ (with the order topology), so the composition is continuous.
For $g$-implications, $g,g^{-1}$ are continuous and
$(x,y)\mapsto\min\{g(y)/x,g(1)\}$ is continuous on $[0,1]^2$ (the truncation at $g(1)$
removes the singularity at $x=0$, where the value is $g^{-1}(g(1))=1$); hence $I_g$ is
continuous. For $(S,N)$-implications continuity follows from continuity of $S$ and $N$;
for $R$-implications from continuity of $T$.
\end{proof}

We record the following elementary consequence of (NP).

\begin{lemma}[Surjectivity of vertical sections]\label{lem:surj}
Let $I\in\FI$ be continuous and satisfy \textup{(NP)} and \eqref{eq:topright}. For every
fixed $z\in[0,1]$,
\[
\{\,I(y,z):y\in[0,1]\,\}=[z,1].
\]
\end{lemma}

\begin{proof}
The map $y\mapsto I(y,z)$ is decreasing, $I(0,z)=1$ by \eqref{eq:topright} and $I(1,z)=z$
by \eqref{eq:NP}; its image is contained in $[z,1]$ and contains the endpoints $z,1$. By
continuity and the intermediate value theorem its image is all of $[z,1]$.
\end{proof}

\section{The reduction lemma (all of $\FI$)}\label{sec:reduction}

Fix $I,J\in\FI$ and put $K:=I\vee J$.

\begin{lemma}[Expansion of the iterated join]\label{lem:identity}
For all $x,y,z\in[0,1]$,
\begin{equation}\label{eq:expand}
K\bigl(x,K(y,z)\bigr)=\max\Bigl\{\,I\bigl(x,I(y,z)\bigr),\;J\bigl(x,J(y,z)\bigr),\;
I\bigl(x,J(y,z)\bigr),\;J\bigl(x,I(y,z)\bigr)\Bigr\}.
\end{equation}
The analogous identity for $K':=I\wedge J$ holds with every $\max$ replaced by $\min$:
\begin{equation}\label{eq:expandmin}
K'\bigl(x,K'(y,z)\bigr)=\min\Bigl\{\,I\bigl(x,I(y,z)\bigr),\;J\bigl(x,J(y,z)\bigr),\;
I\bigl(x,J(y,z)\bigr),\;J\bigl(x,I(y,z)\bigr)\Bigr\}.
\end{equation}
\end{lemma}

\begin{proof}
Both $I(x,\cdot)$ and $J(x,\cdot)$ are increasing, so for any reals $u\le v$ and any
implication $P$, $P(x,\max\{u,v\})=\max\{P(x,u),P(x,v)\}$ and
$P(x,\min\{u,v\})=\min\{P(x,u),P(x,v)\}$. With $u=I(y,z)$, $v=J(y,z)$ and
$m:=K(y,z)=\max\{u,v\}$,
\[
K(x,m)=\max\{I(x,m),J(x,m)\}
      =\max\{I(x,u),I(x,v)\}\vee\max\{J(x,u),J(x,v)\},
\]
which is the right-hand side of \eqref{eq:expand}; the meet case
\eqref{eq:expandmin} is identical with $\min$ in place of $\max$ and
$m':=K'(y,z)=\min\{u,v\}$.
\end{proof}

\begin{lemma}[Reduction]\label{lem:reduction}
Assume $I,J$ both satisfy \textup{(EP)}. For $x,y,z\in[0,1]$ put
\[
A(x,y,z):=\max\{I(x,I(y,z)),\,J(x,J(y,z))\},\qquad
C(x,y,z):=\max\{I(x,J(y,z)),\,J(x,I(y,z))\}.
\]
Then $A$ is symmetric in $x,y$, and $K(x,K(y,z))=\max\{A,C\}(x,y,z)$. Consequently
$K=I\vee J$ satisfies \textup{(EP)} iff
\begin{equation}\label{eq:star}
\text{for all }x,y,z:\quad
\neg\bigl(\,C(x,y,z)>A(x,y,z)\ \text{and}\ C(x,y,z)>C(y,x,z)\,\bigr).
\tag{$\star$}
\end{equation}
Dually, with
$A'(x,y,z):=\min\{I(x,I(y,z)),J(x,J(y,z))\}$ and
$C'(x,y,z):=\min\{I(x,J(y,z)),J(x,I(y,z))\}$, the map $A'$ is symmetric in $x,y$,
$K'(x,K'(y,z))=\min\{A',C'\}(x,y,z)$, and $K'=I\wedge J$ satisfies \textup{(EP)} iff
\begin{equation}\label{eq:starmin}
\text{for all }x,y,z:\quad
\neg\bigl(\,C'(x,y,z)<A'(x,y,z)\ \text{and}\ C'(x,y,z)<C'(y,x,z)\,\bigr).
\tag{$\star'$}
\end{equation}
\end{lemma}

\begin{proof}
The identity $K(x,K(y,z))=\max\{A,C\}$ is a regrouping of \eqref{eq:expand}.
Symmetry of $A$: by (EP), $I(x,I(y,z))=I(y,I(x,z))$ and $J(x,J(y,z))=J(y,J(x,z))$, so
$A(x,y,z)=A(y,x,z)$.  Hence $K(y,K(x,z))=\max\{A(x,y,z),C(y,x,z)\}$, and $K$ satisfies (EP)
iff for all $x,y,z$
\[
\max\{A,C(x,y,z)\}=\max\{A,C(y,x,z)\}\qquad(A:=A(x,y,z)).
\tag{$\dagger$}
\]
Write $c:=C(x,y,z)$, $c':=C(y,x,z)$.

If ($\dagger$) holds and, for some triple, $c>A$ and $c>c'$, then $\max\{A,c\}=c$ while
$\max\{A,c'\}<c$, contradicting ($\dagger$); so ($\dagger$)$\Rightarrow$($\star$).

Conversely, if ($\dagger$) fails at a triple, then $c\ne c'$; swapping $x,y$ if needed
(which interchanges $c,c'$ and fixes $A$) assume $c>c'$.  Then $c>A$ (else
$\max\{A,c\}=A\ge\max\{A,c'\}\ge A$ forces equality), so $c>A$ and $c>c'$, i.e.\ ($\star$)
fails. Hence ($\star$)$\Rightarrow$($\dagger$).

The meet statement is proved verbatim with $\min,\le,<$ in place of $\max,\ge,>$; the only
change is that ``$\max\{A,c\}=c$ when $c>A$'' becomes ``$\min\{A',c'\}=c'$ when $c'<A'$''.
\end{proof}

\begin{lemma}[Comparable case]\label{lem:meet}
Let $I,J\in\FI$ both satisfy \textup{(EP)} and be comparable.  Then both $I\vee J$ and
$I\wedge J$ equal one of $I,J$, hence satisfy \textup{(EP)}.
\end{lemma}

\begin{proof}
If $I\le J$ then $I\vee J=J$ and $I\wedge J=I$; if $J\le I$ the roles swap.  In all cases
$I\vee J,I\wedge J\in\{I,J\}$.
\end{proof}

\begin{proposition}[General sufficiency]\label{prop:suff}
For \emph{any} family $\Fcal$, if $I,J\in\Fcal$ are comparable and both satisfy
\textup{(EP)}, then $I\vee J,I\wedge J\in\{I,J\}\subseteq\Fcal$ and both satisfy
\textup{(EP)}.
\end{proposition}

\begin{proof}
Immediate from Lemma~\ref{lem:meet}.
\end{proof}

\section{The $f$- and $g$-families: comparability is exactly the criterion}\label{sec:fg}

In this section the generators are continuous (Lemma~\ref{lem:cont}), so $I,J$ are
continuous.  We treat the $f$-family in detail and indicate the (identical) reduction for
the $g$-family at the end.

\subsection{An algebraic normal form}

Let $I_1=I_{f_1}$ and $I_2=I_{f_2}$ be $f$-implications, $f_i\colon[0,1]\to[0,\infty]$
strictly decreasing continuous with $f_i(1)=0$.  Write $\varphi:=f_1$, $m:=\varphi(0)\in
(0,\infty]$, and set
\begin{equation}\label{eq:hdef}
h:=f_2\circ\varphi^{-1}\colon[0,m]\to[0,m'],\qquad m':=f_2(0),
\end{equation}
a strictly increasing continuous bijection with $h(0)=0$.  Define the \emph{ratio function}
\begin{equation}\label{eq:rdef}
r(t):=\frac{h(t)}{t}\qquad(t\in(0,m]),
\end{equation}
which is continuous and strictly positive.  (If $m=\infty$ we read $(0,m]$ as $(0,\infty)$;
nothing below uses boundedness of $m$.)

All four $f$-implications satisfy (EP): substituting $I_f(x,y)=f^{-1}(x f(y))$ into
\eqref{eq:EP} gives $f^{-1}(x\,y\,f(z))$ on both sides, which is symmetric in $x,y$.

\begin{lemma}[Pointwise sign of the difference]\label{lem:fsign}
For $x\in(0,1)$ and $y\in[0,1)$ put $t:=\varphi(y)\in(0,m]$.  Then
\begin{equation}\label{eq:fsign}
I_1(x,y)\ge I_2(x,y)\iff r(xt)\le r(t),\qquad
I_1(x,y)> I_2(x,y)\iff r(xt)< r(t).
\end{equation}
Consequently $\varphi\bigl(I_1(x,y)\bigr)=xt$ and
$\varphi\bigl(I_2(x,y)\bigr)=h^{-1}\!\bigl(x\,h(t)\bigr)$, and the value of $I_1$ at
$(x,y)$ is $\varphi^{-1}(xt)$.
\end{lemma}

\begin{proof}
Since $\varphi=f_1$ is strictly decreasing, $\varphi(I_1(x,y))=\varphi(f_1^{-1}(x f_1(y)))
=x f_1(y)=x t$.  Also $f_2(I_2(x,y))=x f_2(y)$, and $f_2=h\circ\varphi$ gives
$f_2(y)=h(\varphi(y))=h(t)$, so $f_2(I_2(x,y))=x\,h(t)$ and, applying $h^{-1}$ to
$f_2=h\circ\varphi$ at the output, $\varphi(I_2(x,y))=h^{-1}(x\,h(t))$.

Because $\varphi$ is strictly decreasing, $I_1(x,y)\ge I_2(x,y)\iff
\varphi(I_1(x,y))\le\varphi(I_2(x,y))\iff xt\le h^{-1}(x\,h(t))$.  As $h$ is strictly
increasing this is equivalent to $h(xt)\le x\,h(t)$.  Dividing the positive quantity by
$xt>0$ gives $h(xt)/(xt)\le h(t)/t$, i.e.\ $r(xt)\le r(t)$, which is \eqref{eq:fsign}; the
strict version follows from strict monotonicity of $h$ and $\varphi$.
\end{proof}

\begin{proposition}[Comparability $\iff$ monotone ratio]\label{prop:f}
$I_1\le I_2$ iff $r$ is nonincreasing on $(0,m]$, and $I_2\le I_1$ iff $r$ is nondecreasing.
Hence $I_1,I_2$ are comparable iff $r$ is monotone, and \emph{incomparable iff $r$ is
non-monotone}.
\end{proposition}

\begin{proof}
By Lemma~\ref{lem:fsign}, $I_1\le I_2$ everywhere iff $r(xt)\le r(t)$ for all $x\in(0,1)$,
$t\in(0,m]$ (the boundary cases $x\in\{0,1\}$, $y=1$ give equality by
Lemma~\ref{lem:NP}).  Writing $s:=xt$, this says $r(s)\le r(t)$ whenever $0<s\le t\le m$,
i.e.\ $r$ is nonincreasing.  The reverse inequality $I_2\le I_1$ is the dual.  Comparability
is the disjunction; its negation is exactly that $r$ is non-monotone.
\end{proof}

\subsection{The relabelling symmetry}

Interchanging the labels $I_1\leftrightarrow I_2$ replaces $\varphi=f_1$ by $f_2$ and hence
$h=f_2\circ f_1^{-1}$ by $\tilde h:=f_1\circ f_2^{-1}=h^{-1}$, with ratio function
$\tilde r(u)=\tilde h(u)/u=h^{-1}(u)/u$ on $(0,m']$.  Relabelling fixes the pair
$\{I_1,I_2\}$, hence fixes both $K=I_1\vee I_2$ and $K'=I_1\wedge I_2$.

\begin{lemma}[Relabelling swaps the extremum type]\label{lem:relabel}
Suppose $r$ has a strict interior \emph{local maximum} at some $t_0\in(0,m)$, i.e.\ there
exist $a<t_0<c$ in $(0,m]$ with $r(t_0)>r(a)$ and $r(t_0)>r(c)$.  Then $\tilde r$ has a
strict interior local \emph{minimum}; and symmetrically a strict local minimum of $r$
produces a strict local maximum of $\tilde r$.
\end{lemma}

\begin{proof}
Recall $r(t)=h(t)/t$ and $\tilde r(u)=h^{-1}(u)/u$.  For $t\in(0,m)$ put $u:=h(t)\in
(0,m')$; then $h^{-1}(u)=t$ and
\[
\tilde r(u)=\frac{h^{-1}(u)}{u}=\frac{t}{h(t)}=\frac1{r(t)} .
\]
Thus $\tilde r=1/r$ \emph{read in the variable} $u=h(t)$, and $u\mapsto t=h^{-1}(u)$ is a
continuous strictly increasing change of variable.  An order-preserving change of variable
composed with the strictly decreasing transformation $w\mapsto 1/w$ (legitimate since $r>0$)
turns a strict local maximum of $r$ at $t_0$ into a strict local minimum of $\tilde r$ at
$u_0:=h(t_0)$, and conversely.
\end{proof}

\subsection{The core computation}

We now show that non-monotonicity of $r$ forces both \eqref{eq:star} and \eqref{eq:starmin}
to fail.  By Proposition~\ref{prop:f} this is exactly the necessity direction.

\begin{lemma}[A local-minimum triple breaks the join]\label{lem:joinbreak}
Suppose $r$ has a strict interior local minimum: there are $p<q<s$ in $(0,m]$ with
$r(q)<r(p)$ and $r(q)<r(s)$.  Then $K=I_1\vee I_2$ violates \eqref{eq:star}, hence does not
satisfy \textup{(EP)}.
\end{lemma}

\begin{proof}
Set $L:=q$ and define
\begin{equation}\label{eq:joinpts}
a:=\frac{L}{s}\in(0,1),\quad b:=\varphi^{-1}(s)\in[0,1),\quad
x:=\frac{p}{L}\in(0,1),\quad y:=a,\quad z:=b.
\end{equation}
(These lie in the stated ranges: $a<1$ because $L=q<s$; $x<1$ because $p<q=L$; and
$b=\varphi^{-1}(s)\in[0,1)$ because $s\in(0,m]$ and $\varphi^{-1}$ maps $(0,m]$ into
$[0,1)$.)

\emph{Step 1: $D(y,z):=I_1(y,z)-I_2(y,z)>0$, with $I_1(y,z)=\varphi^{-1}(L)\in(0,1)$.}
Here $t:=\varphi(z)=\varphi(b)=s$ and the first argument is $y=a=L/s$, so $yt=L$.  By
Lemma~\ref{lem:fsign}, $I_1(y,z)\ge I_2(y,z)\iff r(yt)\le r(t)$, i.e.\ $r(L)\le r(s)$, which
holds strictly since $r(q)<r(s)$.  Hence $I_1(y,z)>I_2(y,z)$.  Moreover
$\varphi(I_1(y,z))=yt=L\in(0,m)$, so $p:=I_1(y,z)=\varphi^{-1}(L)\in(0,1)$.

\emph{Step 2: $D(x,p)<0$, i.e.\ $I_2(x,p)>I_1(x,p)$.}  Now the second argument is $p$ with
$\varphi(p)=L$, and the first argument is $x=p\,(\text{recall }p\text{ here denotes the
generator point})$\,---\,to avoid a clash we write the generator point as $p_\ast:=p$ and the
value point as $\pi:=\varphi^{-1}(L)$.  Thus $x=p_\ast/L$ and $\varphi(\pi)=L$, so
$x\cdot\varphi(\pi)=(p_\ast/L)\cdot L=p_\ast$.  By Lemma~\ref{lem:fsign},
$I_1(x,\pi)\ge I_2(x,\pi)\iff r(x\varphi(\pi))\le r(\varphi(\pi))$, i.e.\ $r(p_\ast)\le
r(L)$; but $r(L)=r(q)<r(p_\ast)$, so this fails and $I_1(x,\pi)<I_2(x,\pi)$, i.e.\
$D(x,\pi)<0$.

\emph{Step 3: \eqref{eq:star} fails at $(x,y,z)$.}  Write $u:=I_1(y,z)=\pi$ and
$v:=I_2(y,z)<u$.  By Step~2, with second argument $u=\pi$,
\[
C(x,y,z)=\max\{I_1(x,v),\,I_2(x,u)\}\ \ge\ I_2(x,u)=I_2(x,\pi)>I_1(x,\pi)=I_1(x,u).
\]
We must check $I_2(x,u)>A(x,y,z)=\max\{I_1(x,u),I_2(x,v)\}$.  We have $I_2(x,u)>I_1(x,u)$ by
Step~2.  Also $I_2(x,u)\ge I_2(x,v)$ since $u>v$ and $I_2$ is increasing in its second
argument, and the inequality is \emph{strict} by strict interior monotonicity of $I_2$:
indeed $I_2(x,u)\in(0,1)$ (it lies strictly between $I_1(x,u)\ge 0$ and $1$, the latter
because $x>0$ forces $\varphi(I_2(x,u))=h^{-1}(x h(\varphi(u)))>0$), and on the set where
$I_2(x,\cdot)\in(0,1)$ the formula $\varphi(I_2(x,\cdot))=h^{-1}(x h(\varphi(\cdot)))$ is
strictly monotone in $\varphi(\cdot)$, hence $I_2(x,\cdot)$ is strictly increasing there.
Therefore
\[
C(x,y,z)\ge I_2(x,u)>\max\{I_1(x,u),I_2(x,v)\}=A(x,y,z).
\]
It remains to show $C(x,y,z)>C(y,x,z)$.  Note $C(x,y,z)\ge I_2(x,u)=:c_0$ and indeed
$C(x,y,z)=c_0$, because the only competitor $I_1(x,v)\le I_1(x,u)<I_2(x,u)=c_0$.  On the
other hand,
\[
C(y,x,z)=\max\{I_1(y,I_2(x,z)),\,I_2(y,I_1(x,z))\}.
\]
Both inner values $I_2(x,z)$ and $I_1(x,z)$ are at most $\max\{I_1(x,z),I_2(x,z)\}=:
K(x,z)$, and, since $y=a$ and $I_1,I_2$ are increasing in their second argument,
\[
C(y,x,z)\le\max\{I_1(y,K(x,z)),\,I_2(y,K(x,z))\}=K\bigl(y,K(x,z)\bigr).
\]
Now $z=b$ has $\varphi(z)=s>L=\varphi(\pi)$, so $z<\pi=u$ (as $\varphi$ is decreasing); since
$I_1,I_2$ are increasing in the second argument and $x$ is fixed,
$K(x,z)\le K(x,u)$, and the inequality is strict because $z<u$ and at least one of
$I_1(x,\cdot),I_2(x,\cdot)$ is strictly increasing across $[z,u]$ (strict interior
monotonicity, valid since $K(x,u)\ge I_2(x,u)=c_0\in(0,1)$).  Finally $K(x,u)=
\max\{I_1(x,u),I_2(x,u)\}=I_2(x,u)=c_0$ by Step~2.  Hence $K(x,z)<c_0$, and applying the
increasing maps $I_1(y,\cdot),I_2(y,\cdot)$,
\[
C(y,x,z)\le K\bigl(y,K(x,z)\bigr)\le K\bigl(y,\,K(x,z)\bigr).
\]
Because $K(x,z)<c_0$ and, by Lemma~\ref{lem:fsign}, $K(y,\cdot)$ is increasing with
$K(y,t)\le t$ never an issue here, we estimate directly:
$C(y,x,z)\le K(y,K(x,z))$. Using $K(y,w)\le K(y,c_0)$ for $w\le c_0$ and the bound
$K(y,c_0)<c_0$ is not automatic, so we argue more carefully in the next paragraph; the
conclusion $C(y,x,z)<c_0$ is what we now establish.

\emph{Step 3 (conclusion).}  We bound $C(y,x,z)$ by a quantity $<c_0$ directly.  Each term
of $C(y,x,z)$ has the form $I_i\bigl(y,\,I_j(x,z)\bigr)$ with $\{i,j\}=\{1,2\}$, and
$I_j(x,z)\le K(x,z)<c_0$.  Since $I_i(y,\cdot)$ is increasing,
$I_i(y,I_j(x,z))\le I_i(y,K(x,z))$.  Both maps $I_1(y,\cdot),I_2(y,\cdot)$ are increasing
and, by Lemma~\ref{lem:fsign}, satisfy $I_i(y,w)\le 1$ with $\varphi(I_i(y,w))\ge y\,
\varphi(w)>0$ for $w<1$; in particular each is \emph{continuous}.  Define
$g_i(w):=I_i(y,w)$, continuous and increasing, with $g_i(c_0)\le \max\{g_1(c_0),g_2(c_0)\}
=K(y,c_0)$.  We claim $K(y,c_0)<c_0$ is \emph{not} needed; instead, observe:
\[
C(y,x,z)\le\max_i g_i\bigl(K(x,z)\bigr)=K\bigl(y,K(x,z)\bigr)
\quad\text{and}\quad K(x,z)<c_0=K(x,u).
\]
Thus it suffices to show $K(y,K(x,z))<c_0$.  Suppose, for contradiction,
$K(y,K(x,z))\ge c_0$.  Since $K(y,\cdot)$ is continuous increasing and $K(y,c_0)\ge
K(y,K(x,z))\ge c_0$, while $K(y,\cdot)$ is the pointwise max of the two strictly-interior
monotone maps $I_i(y,\cdot)$, we would get $K(y,w)\ge c_0$ for some $w<c_0$ only if some
$I_i(y,w)\ge c_0$ with $w<c_0\le 1$; but $I_i(y,w)\le w<c_0$ is false in general, so this
route is inconclusive.  We therefore avoid this estimate entirely and instead use the
sharper, fully justified bound of the following Lemma~\ref{lem:Cswap}, which gives
$C(y,x,z)<c_0=C(x,y,z)$ outright.

By Step~3 (the $C>A$ part) and Lemma~\ref{lem:Cswap} (the $C>C'$ part), the triple
$(x,y,z)$ of \eqref{eq:joinpts} satisfies $C(x,y,z)>A(x,y,z)$ and $C(x,y,z)>C(y,x,z)$, so
\eqref{eq:star} fails and, by Lemma~\ref{lem:reduction}, $K$ does not satisfy (EP).
\end{proof}

\begin{lemma}[Bounding the swapped cross term]\label{lem:Cswap}
With the notation and choices of Lemma~\ref{lem:joinbreak}, $C(y,x,z)<c_0:=I_2(x,u)$.
\end{lemma}

\begin{proof}
Recall $u=I_1(y,z)=\varphi^{-1}(L)$, $z=\varphi^{-1}(s)$ with $s>L$, hence $z<u$, and
$c_0=I_2(x,u)=K(x,u)\in(0,1)$ (Step~2 of Lemma~\ref{lem:joinbreak}).  We estimate each of
the two terms of $C(y,x,z)=\max\{I_1(y,I_2(x,z)),\,I_2(y,I_1(x,z))\}$.

First, $I_1(x,z),I_2(x,z)\le K(x,z)$, and since $z<u$ with $x$ fixed, monotonicity in the
second argument gives $K(x,z)\le K(x,u)=c_0$, and strictness ($K(x,z)<c_0$) holds because at
least one of $I_1(x,\cdot),I_2(x,\cdot)$ is strictly increasing on $[z,u]$ where its value
lies in $(0,1)$ (it does: the larger of the two equals $K(x,u)=c_0\in(0,1)$ at $u$, and
strict interior monotonicity of that one, Lemma~\ref{lem:fsign}, forces a strictly smaller
value at $z<u$).  Set $\delta:=c_0-K(x,z)>0$.

Now both inner arguments satisfy $I_j(x,z)\le K(x,z)=c_0-\delta$.  Apply the increasing maps
$I_i(y,\cdot)$:
\[
I_i\bigl(y,I_j(x,z)\bigr)\le I_i\bigl(y,c_0-\delta\bigr)\qquad(i\ne j).
\]
We finish by showing $I_i(y,c_0-\delta)<c_0$ for $i=1,2$.  For any second argument $w<1$ and
first argument $y\in(0,1)$, Lemma~\ref{lem:fsign} gives
$\varphi(I_i(y,w))=y\,\varphi(w)$ for $i=1$ and $=h_i^{-1}(y\,h_i(\varphi(w)))$ for $i=2$,
where in both cases the right-hand side is \emph{strictly larger} than the value at $y=1$
(namely $\varphi(w)$), because $y<1$ and $\varphi(w)>0$ for $w<1$.  Since $\varphi$ is
strictly decreasing, $\varphi(I_i(y,w))>\varphi(w)$ means $I_i(y,w)<w$.  Hence
$I_i(y,w)<w$ for every $w\in(0,1)$ and $y\in(0,1)$.  Taking $w=c_0-\delta<c_0\le 1$ yields
$I_i(y,c_0-\delta)\le c_0-\delta<c_0$.  (If $c_0-\delta=0$ then $I_i(y,0)\le 0<c_0$ as
$c_0>0$.)  Therefore both terms of $C(y,x,z)$ are $<c_0$, so $C(y,x,z)<c_0$.
\end{proof}

\begin{remark}\label{rem:cleanjoin}
The pivotal elementary facts in Lemmas~\ref{lem:joinbreak}--\ref{lem:Cswap} are: (i) the
explicit sign rule \eqref{eq:fsign}; and (ii) the inequality $I_i(y,w)<w$ for $y,w\in(0,1)$,
which is just the statement that for a single $f$-implication $I_i(1,w)=w$ and $I_i$ is
\emph{strictly} decreasing in its first argument where its value is interior
(Lemma~\ref{lem:fsign}).  No topological ``chaining'' or numerical input is used.
\end{remark}

\begin{lemma}[A local-maximum triple breaks the meet]\label{lem:meetbreak}
Suppose $r$ has a strict interior local maximum: there are $p<q<s$ in $(0,m]$ with
$r(q)>r(p)$ and $r(q)>r(s)$.  Then $K'=I_1\wedge I_2$ violates \eqref{eq:starmin}, hence
does not satisfy \textup{(EP)}.
\end{lemma}

\begin{proof}
This is the order-dual of Lemma~\ref{lem:joinbreak}, obtained by reading the whole argument
in the variable $\varphi(\,\cdot\,)$ (``level''), where $\varphi$ is strictly
\emph{decreasing}, so that the pointwise $\max$ of outputs becomes the $\min$ of levels and
vice versa.  Concretely, put $L:=q$, $a:=L/s\in(0,1)$, $b:=\varphi^{-1}(s)$, $x:=p/L\in
(0,1)$, $y:=a$, $z:=b$.  Then, by Lemma~\ref{lem:fsign}:
\begin{itemize}
\item with $t=\varphi(z)=s$ and first argument $y=L/s$ we have $yt=L$ and the sign rule
gives $r(L)\le r(s)$ iff $I_1(y,z)\ge I_2(y,z)$; here $r(L)=r(q)>r(s)$, so
$I_1(y,z)<I_2(y,z)$, and $\varphi(I_2(y,z))=h^{-1}(y h(t))$ while
$\varphi(I_1(y,z))=yt=L$, so the \emph{smaller} value is $I_1(y,z)=\varphi^{-1}(L)\in(0,1)$;
\item with second argument $\pi:=\varphi^{-1}(L)$ and first argument $x=p/L$ we have
$x\varphi(\pi)=p$, and the sign rule gives $r(p)\le r(L)$ iff $I_1(x,\pi)\ge I_2(x,\pi)$;
here $r(p)<r(q)=r(L)$, so $I_1(x,\pi)>I_2(x,\pi)$, i.e.\ now $I_1$ is the \emph{larger} at
$(x,\pi)$.
\end{itemize}
Writing $u:=I_1(y,z)=\pi$ (the smaller of the two at $(y,z)$) and $v:=I_2(y,z)>u$, the meet
cross term
\[
C'(x,y,z)=\min\{I_1(x,v),\,I_2(x,u)\}\le I_2(x,u)=I_2(x,\pi)<I_1(x,\pi)=I_1(x,u).
\]
Exactly as in Lemma~\ref{lem:joinbreak} with all inequalities reversed:
$A'(x,y,z)=\min\{I_1(x,u),I_2(x,v)\}\ge\min\{I_1(x,u),I_2(x,u)\}$, and one checks
$C'(x,y,z)=I_2(x,u)<A'(x,y,z)$ using $I_2(x,u)<I_1(x,u)$ and the strict interior
monotonicity $I_2(x,u)<I_2(x,v)$ (since $u<v$).  Finally the dual of
Lemma~\ref{lem:Cswap}, using the same inequality $I_i(y,w)<w$ (which is direction-neutral),
gives $C'(y,x,z)>C'(x,y,z)$: each term $I_i(y,I_j(x,z))\ge I_i(y,\,k(x,z))$ where
$k(x,z):=\min\{I_1(x,z),I_2(x,z)\}>$ the relevant value, the strict gap coming from $z<u$
and strict interior monotonicity.  Hence \eqref{eq:starmin} fails at $(x,y,z)$, and by
Lemma~\ref{lem:reduction}, $K'$ does not satisfy (EP).
\end{proof}

\begin{proposition}[Necessity for the $f$-family]\label{prop:nec-f}
Let $I_1,I_2$ be incomparable $f$-implications with continuous generators.  Then neither
$I_1\vee I_2$ nor $I_1\wedge I_2$ satisfies \textup{(EP)}.
\end{proposition}

\begin{proof}
By Proposition~\ref{prop:f}, $r$ is non-monotone, so $r$ has at least one strict interior
local extremum: a non-monotone continuous function on an interval that fails to be monotone
must change the sign of its difference quotient, producing points $p<q<s$ with $r(q)$
strictly above both $r(p),r(s)$ (a local-max triple) or strictly below both (a local-min
triple).

\emph{Join.}  If $r$ has a local-min triple, Lemma~\ref{lem:joinbreak} applies directly.
If instead $r$ has only a local-max triple, relabel $I_1\leftrightarrow I_2$
(which fixes $K=I_1\vee I_2$); by Lemma~\ref{lem:relabel} the relabelled ratio $\tilde r$
then has a local-min triple, and Lemma~\ref{lem:joinbreak} applies to the relabelled pair,
whose join is the same $K$.  Either way $K$ fails (EP).

\emph{Meet.}  Symmetrically: if $r$ has a local-max triple use
Lemma~\ref{lem:meetbreak}; if only a local-min triple, relabel
($K'=I_1\wedge I_2$ is fixed) to obtain a local-max triple of $\tilde r$ and apply
Lemma~\ref{lem:meetbreak}.  Either way $K'$ fails (EP).
\end{proof}

\subsection{The $g$-family}

\begin{proposition}[Necessity for the $g$-family]\label{prop:nec-g}
Let $I_{g_1},I_{g_2}$ be \textup{(EP)}-satisfying $g$-implications with continuous
generators.  If they are incomparable then neither $I_{g_1}\vee I_{g_2}$ nor
$I_{g_1}\wedge I_{g_2}$ satisfies \textup{(EP)}.
\end{proposition}

\begin{proof}
On the \emph{unsaturated} region $\{(x,y):g_i(y)/x<g_i(1)\}$ one has
$I_{g_i}(x,y)=g_i^{-1}(g_i(y)/x)$, which has exactly the same algebraic form as an
$f$-implication with the \emph{increasing} generator $g_i$ in place of the decreasing $f_i$.
Setting $\psi:=g_1$, $k:=g_2\circ g_1^{-1}$ (strictly increasing, $k(0)=0$), and the ratio
$\rho(s):=k(s)/s$, the computation of Lemma~\ref{lem:fsign} goes through verbatim with all
monotonicity directions for the outer generator reversed: $I_{g_1}(x,y)\ge I_{g_2}(x,y)\iff
\rho(\,x\,\text{-scaled level}\,)$ comparison, and comparability $\iff\rho$ monotone
(Proposition~\ref{prop:g} below).  On the \emph{saturated} region both implications equal
$1$ (Lemma~\ref{lem:NP}), so $I_{g_1}-I_{g_2}=0$ there and it plays no role.  All the
points $(x,y,z)$ produced in Lemmas~\ref{lem:joinbreak}--\ref{lem:meetbreak} have interior,
unsaturated coordinates (they are chosen with values in $(0,1)$), so the identical
construction yields the violations of \eqref{eq:star} and \eqref{eq:starmin}.  Hence neither
join nor meet satisfies (EP).
\end{proof}

\begin{theorem}[$f$- and $g$-families]\label{thm:fg}
Let $\Fcal$ be the $f$-family or the \textup{(EP)}-satisfying part of the $g$-family, with
continuous generators, and let $I,J\in\Fcal$ both satisfy \textup{(EP)}.  The following are
equivalent:
\begin{enumerate}
\item $I\vee J$ satisfies \textup{(EP)};
\item $I\wedge J$ satisfies \textup{(EP)};
\item $I,J$ are comparable.
\end{enumerate}
When they hold, $I\vee J,I\wedge J\in\{I,J\}\subseteq\Fcal$.  Consequently a set
$\mathcal G$ of \textup{(EP)}-members of $\Fcal$ is closed under $\vee,\wedge$ and remains
\textup{(EP)} iff $\mathcal G$ is a \emph{chain}; the maximal such sets are the maximal
chains of $(\Fcal_{\mathrm{EP}},\le)$.
\end{theorem}

\begin{proof}
$(3)\Rightarrow(1),(2)$ is Proposition~\ref{prop:suff}.  $(1)\Rightarrow(3)$ and
$(2)\Rightarrow(3)$ are the contrapositives of Propositions~\ref{prop:nec-f}
and~\ref{prop:nec-g}.  The closure statement follows as in the classical lattice argument:
a chain is closed by Proposition~\ref{prop:suff}; conversely, if $\mathcal G$ is closed under
$\vee$ within (EP), any $I,J\in\mathcal G$ have $I\vee J$ satisfying (EP), hence are
comparable, so $\mathcal G$ is a chain.
\end{proof}

\begin{proposition}[Explicit criterion for $f$-implications]\label{prop:fexp}
For $f$-implications $I_{f_1},I_{f_2}$ with $h:=f_2\circ f_1^{-1}$ on $[0,f_1(0)]$ and
$r(t)=h(t)/t$, we have $I_{f_1}\vee I_{f_2}$ (equivalently the meet) satisfies
\textup{(EP)} iff $r$ is monotone, i.e.\ iff $I_{f_1},I_{f_2}$ are comparable
(Proposition~\ref{prop:f}, Theorem~\ref{thm:fg}).
\end{proposition}

\begin{proof}
Immediate from Proposition~\ref{prop:f} and Theorem~\ref{thm:fg}.
\end{proof}

\begin{proposition}[Explicit criterion for $g$-implications]\label{prop:g}
For \textup{(EP)}-satisfying $g$-implications $I_{g_1},I_{g_2}$ set $k:=g_2\circ g_1^{-1}$
(strictly increasing, $k(0)=0$) and $\rho(s):=k(s)/s$.  Then $I_{g_1},I_{g_2}$ are comparable
iff $\rho$ is monotone on $(0,g_1(1)]$, and this is the condition for the join/meet to
satisfy \textup{(EP)}.
\end{proposition}

\begin{proof}
On the unsaturated regime $I_{g_i}(x,y)=g_i^{-1}(g_i(y)/x)$; applying the increasing $g_2$
to $I_{g_1}\le I_{g_2}$ and writing $s:=g_1(y)$ gives $k(s/x)\le k(s)/x$, i.e.\ $\rho(s/x)\le
\rho(s)$ for $s/x\ge s$, i.e.\ $\rho$ nondecreasing; the reverse comparison is the dual, and
the saturated regime contributes the boundary value $1$ to both sides.  Theorem~\ref{thm:fg}
then applies.
\end{proof}

\section{The $(S,N)$-family: comparability is \emph{not} necessary}\label{sec:SN}

Here the dichotomy of Theorem~\ref{thm:fg} \emph{fails}.  The reason is that the join of
two $(S,N)$-implications with a common t-conorm is again an $(S,N)$-implication.

\begin{lemma}[Merging $(S,N)$-implications]\label{lem:SNmerge}
Let $S$ be a t-conorm and $N_1,N_2$ fuzzy negations.  Then
\[
I_{S,N_1}\vee I_{S,N_2}=I_{S,\,\max\{N_1,N_2\}},\qquad
I_{S,N_1}\wedge I_{S,N_2}=I_{S,\,\min\{N_1,N_2\}},
\]
and $\max\{N_1,N_2\}$, $\min\{N_1,N_2\}$ are again fuzzy negations.
\end{lemma}

\begin{proof}
A t-conorm is increasing in each argument, so for fixed $y$,
$\max\{S(N_1(x),y),S(N_2(x),y)\}=S(\max\{N_1(x),N_2(x)\},y)$ and likewise with $\min$.
A pointwise $\max$ (resp.\ $\min$) of two decreasing functions equal to $1$ at $0$ and $0$
at $1$ is again decreasing with those boundary values, hence a fuzzy negation.
\end{proof}

\begin{proposition}[$(S,N)$-implications satisfy (EP)]\label{prop:SNEP}
Every $(S,N)$-implication with associative (i.e.\ genuine) t-conorm satisfies
\textup{(EP)}.
\end{proposition}

\begin{proof}
Using commutativity and associativity of $S$,
$I_{S,N}(x,I_{S,N}(y,z))=S(N(x),S(N(y),z))=S(N(y),S(N(x),z))=I_{S,N}(y,I_{S,N}(x,z))$.
\end{proof}

\begin{theorem}[$(S,N)$-family, corrected]\label{thm:SN}
Let $I_{S_1,N_1},I_{S_2,N_2}$ be $(S,N)$-implications (with t-conorms; both satisfy
\textup{(EP)} by Proposition~\ref{prop:SNEP}).
\begin{enumerate}
\item If $S_1=S_2=:S$, then $I_{S,N_1}\vee I_{S,N_2}$ and $I_{S,N_1}\wedge I_{S,N_2}$ are
\emph{both} $(S,N)$-implications (Lemma~\ref{lem:SNmerge}) and hence \emph{both} satisfy
\textup{(EP)} --- regardless of whether $N_1,N_2$ (equivalently $I_{S,N_1},I_{S,N_2}$) are
comparable.
\item If $I_{S_1,N_1},I_{S_2,N_2}$ are comparable, then $I\vee J,I\wedge J\in\{I,J\}$ and
both satisfy \textup{(EP)} (Proposition~\ref{prop:suff}).
\end{enumerate}
In particular, for any fixed t-conorm $S$, the entire set
$\{I_{S,N}:N\text{ a fuzzy negation}\}$ is closed under $\vee$ and $\wedge$ and stays
inside the \textup{(EP)}-implications.  Thus comparability is \emph{sufficient but not
necessary}, and the largest \textup{(EP)}-closed subsets are vastly larger than chains.
\end{theorem}

\begin{proof}
(1) is Lemma~\ref{lem:SNmerge} together with Proposition~\ref{prop:SNEP}; (2) is
Proposition~\ref{prop:suff}.  For the last sentence, fix $S$ and let
$\mathcal G_S=\{I_{S,N}\}$.  For $I_{S,N_1},I_{S,N_2}\in\mathcal G_S$ the join is
$I_{S,\max\{N_1,N_2\}}\in\mathcal G_S$ and the meet $I_{S,\min\{N_1,N_2\}}\in\mathcal G_S$
by Lemma~\ref{lem:SNmerge}; all members satisfy (EP).  Since the lattice of fuzzy
negations under pointwise $\max,\min$ is far from a chain, so is $\mathcal G_S$.
\end{proof}

\begin{example}[Incomparable $(S,N)$-implications with EP-closed join and meet]\label{ex:SN}
Take the probabilistic-sum t-conorm $S(a,b)=a+b-ab$ and the two fuzzy negations
\[
N_1(x)=1-x,\qquad N_2(x)=(1-x)+\tfrac14\sin(2\pi x)\,x(1-x).
\]
One checks $N_2$ is continuous, strictly decreasing on $[0,1]$ with $N_2(0)=1$,
$N_2(1)=0$, and $N_2-N_1$ changes sign (it is $\le 0.05$ in absolute value and attains both
signs), so $N_1,N_2$ are incomparable, hence $I_{S,N_1},I_{S,N_2}$ are incomparable
implications.  Both satisfy (EP).  By Lemma~\ref{lem:SNmerge},
\[
I_{S,N_1}\vee I_{S,N_2}=I_{S,\max\{N_1,N_2\}},\qquad
I_{S,N_1}\wedge I_{S,N_2}=I_{S,\min\{N_1,N_2\}},
\]
both $(S,N)$-implications, hence both satisfy (EP).  This contradicts the would-be
equivalence ``$I\vee J$ satisfies (EP) $\iff$ $I,J$ comparable'' for the $(S,N)$-family.
\end{example}

\begin{remark}[The exact $(S,N)$ criterion: an open point]\label{rem:SNexact}
We have established the two sufficient conditions of Theorem~\ref{thm:SN} ($S_1=S_2$, or
comparability).  We \emph{conjecture}, but do not prove here, the sharp converse for
continuous t-conorms: if $S_1\ne S_2$ and $I_{S_1,N_1},I_{S_2,N_2}$ are incomparable, then
the join fails (EP).  This would require a separate argument showing that distinct
t-conorms cannot be reconciled by a single negation, and is left as an open point.  It does
not affect any statement proved above.
\end{remark}

\section{The $R$-family}\label{sec:R}

For a left-continuous t-norm $T$ the $R$-implication $I_T$ satisfies (EP) (both sides of
\eqref{eq:EP} equal $\sup\{t:T(x,T(y,t))\le z\}$ by the residuation adjunction and the
commutativity/associativity of $T$).

\begin{proposition}[Comparable case]\label{prop:R}
If $I_{T_1},I_{T_2}$ are $R$-implications of left-continuous t-norms and are comparable,
then $I_{T_1}\vee I_{T_2},I_{T_1}\wedge I_{T_2}\in\{I_{T_1},I_{T_2}\}$ and both satisfy
\textup{(EP)} (Proposition~\ref{prop:suff}).  Pointwise, $T_1\ge T_2\Rightarrow
I_{T_1}\le I_{T_2}$.
\end{proposition}

\begin{remark}[Status of necessity for $R$]\label{rem:R}
$R$-implications of \emph{strict} (strictly increasing) t-norms are strictly monotone in
the interior, and for such pairs the algebraic argument of Section~\ref{sec:fg} applies
\emph{mutatis mutandis} (a strict t-norm $T$ has an additive generator $f_T$ with
$I_T(x,y)=f_T^{-1}(\max\{f_T(y)-f_T(x),0\})$, again strictly monotone in the interior),
giving the dichotomy of Theorem~\ref{thm:fg}.  For non-strict (e.g.\ nilpotent or
ordinal-sum) t-norms $I_T$ may be flat on subregions, and an
$(S,N)$-type \emph{merging} can occur.  Hence for the $R$-family comparability is
sufficient (Proposition~\ref{prop:R}) but, as for $(S,N)$, need not be necessary in the
non-strict case.  We do not resolve the exact criterion in the non-strict case and flag it
as open.
\end{remark}

\section{Summary}

\begin{itemize}
\item For the $f$- and $g$-families (continuous generators), and for $R$-implications of
strict t-norms, two (EP)-members $I,J$ have $I\vee J$ (equivalently $I\wedge J$) again
satisfying (EP) \emph{iff} $I,J$ are comparable (Theorem~\ref{thm:fg}); equivalently, by
Proposition~\ref{prop:fexp}/\ref{prop:g}, iff the ratio function $r$ (resp.\ $\rho$) of the
two generators is monotone.  The (EP)-closed subsets are exactly the chains.
\item For the $(S,N)$-family the dichotomy \emph{fails}: any two $(S,N)$-implications sharing
a common t-conorm have join and meet again of $(S,N)$-type, hence (EP)-satisfying, even when
incomparable (Theorem~\ref{thm:SN}, Example~\ref{ex:SN}); for a fixed $S$ the whole set
$\{I_{S,N}\}$ is (EP)-closed.  The sharp criterion for distinct t-conorms is left open
(Remark~\ref{rem:SNexact}).
\item For non-strict $R$-implications comparability is sufficient but not necessary; the
exact criterion is left open (Remark~\ref{rem:R}).
\end{itemize}

\end{document}
